\unitopen{40}
\begin{candos}
\item odróżnić stronę bierną czynności (\de{Vorgangspassiv}) od strony biernej stanu (\de{Zustandspassiv});
\item użyć strony biernej z czasownikami modalnymi (\de{muss geprüft werden});
\item zastąpić stronę bierną konstrukcjami zastępczymi (\de{sich lassen}, \de{-bar}, \de{man});
\item opisać wymagania i stan systemu w piśmie formalnym.
\end{candos}

\section{Słownictwo}
\vocabtable{40}

\section{Czynność kontra stan}
\textbf{\de{Vorgangspassiv}} (\de{werden} + PII) opisuje \emph{proces}: \de{Die Tür wird geschlossen.} \textbf{\de{Zustandspassiv}} (\de{sein} + PII) opisuje
\emph{stan po procesie}: \de{Die Tür ist geschlossen.}

\begin{center}\small
\begin{tabular}{@{}lll@{}}
\toprule
\textsf{\textbf{Czas}} & \textsf{\textbf{Vorgangspassiv}} & \textsf{\textbf{Zustandspassiv}}\\
\midrule
Präsens & \de{Der Test wird ausgeführt.} & \de{Der Test ist ausgeführt.}\\
Präteritum & \de{Der Test wurde ausgeführt.} & \de{Der Test war ausgeführt.}\\
Perfekt & \de{Der Test ist ausgeführt worden.} & \de{Der Test ist ausgeführt gewesen.} (rzadko)\\
\bottomrule
\end{tabular}
\end{center}
\textbf{Z modalnym:} \de{Der Fehler muss behoben werden.} \de{Das Projekt kann nicht abgeschlossen werden.} \de{Die Daten dürfen nicht
gelöscht werden.} Perfekt: \de{Der Fehler hat behoben werden müssen.}

\section{Zastępniki strony biernej}
\begin{center}\small
\begin{tabular}{@{}ll@{}}
\toprule
\textsf{\textbf{Konstrukcja}} & \textsf{\textbf{Przykład}}\\
\midrule
\de{man} & \de{Man kann das lösen.}\\
\de{sich lassen} & \de{Das lässt sich lösen.} (= \de{Das kann gelöst werden.})\\
\de{-bar}, \de{-lich} & \de{lösbar}, \de{machbar}, \de{vermeidbar}\\
\de{sein + zu + Inf.} & \de{Der Fehler ist zu beheben.} (= musi być usunięty)\\
\de{es gibt + zu + Inf.} & \de{Es gibt viel zu tun.}\\
\bottomrule
\end{tabular}
\end{center}

\begin{zasada}[Wybór]
Wskazujesz \emph{kto} to zrobił: strona czynna. Ważny jest \emph{proces}: \de{werden}-Passiv. Ważny jest \emph{stan}: \de{sein}-Passiv.
Chcesz zaznaczyć możliwość: \de{sich lassen} lub \de{-bar}. Uważaj, aby nie mieszać: \de{Die Tür ist geschlossen worden} znaczy
„została zamknięta” (czynność w przeszłości), \de{Die Tür ist geschlossen} znaczy „jest zamknięta”.
\end{zasada}

\begin{falle}
Polski nie rozróżnia ich tak wyraźnie („drzwi są zamknięte” vs „drzwi zostały zamknięte”): w niemieckim to dwie różne konstrukcje. Błąd:
\de{Das Büro wird geschlossen} jako stan (znaczy „jest zamykane”). Drugi: \de{Der Fehler kann behoben} (brak \de{werden}).
Trzeci: \de{zu}-konstrukcja nie zastępuje \de{werden}-Passiv w czasie przeszłym.
\end{falle}

\begin{szwajcaria}
Szwajcarskie instrukcje, regulaminy i ogłoszenia używają bardzo często konstrukcji z \de{ist \dots{} zu}: \de{Der Eingang ist freizuhalten.}
(wejście należy pozostawić wolne). Znajomość tych form skraca czas czytania umów i regulaminów budynku.
\end{szwajcaria}

\begin{praca}
\de{Die Konfiguration ist abgeschlossen.} (stan) \de{Die Konfiguration wird gerade abgeschlossen.} (proces) \de{Das Problem lässt sich
nicht reproduzieren.} \de{Die Dokumentation ist bis Freitag zu aktualisieren.} \de{Die Daten dürfen nicht ohne Genehmigung gelöscht werden.}
\end{praca}

\section{Dialogi}
\textbf{Dialog 1 --- status projektu}
\begin{dialog}
\turn{A}{Ist die Migration schon abgeschlossen?}{Czy migracja jest już zakończona?}
\turn{B}{Nein, sie wird noch durchgeführt. Die Daten müssen erst geprüft werden.}{Nie, jest jeszcze przeprowadzana. Dane muszą najpierw zostać sprawdzone.}
\turn{A}{Und der Fehler im Login?}{A błąd w logowaniu?}
\turn{B}{Der ist behoben. Das hat sich schnell lösen lassen.}{Jest naprawiony. Udało się to szybko rozwiązać.}
\end{dialog}

\section{Ćwiczenia}
\begin{cwiczenia}
\exitem{Zustandspassiv: \de{Die Tür wird geschlossen.} (stan)}{\de{Die Tür ist geschlossen.}}
\exitem{Vorgangspassiv z modalnym: \de{Man muss den Fehler beheben.}}{\de{Der Fehler muss behoben werden.}}
\exitem{Zastępnik: \de{Man kann das lösen.} z \de{sich lassen}.}{\de{Das lässt sich lösen.}}
\exitem{Zastępnik: \de{Der Fehler muss behoben werden.} z \de{sein + zu}.}{\de{Der Fehler ist zu beheben.}}
\exitem{Utwórz przymiotnik z \de{-bar}: \de{lösen}.}{\de{lösbar}}
\exitem{Co znaczy: \de{Das Büro ist geschlossen.}}{Biuro jest zamknięte (stan).}
\exitem{Co znaczy: \de{Das Büro wird geschlossen.}}{Biuro jest zamykane (proces).}
\exitem{Popraw: \de{Der Fehler kann behoben.}}{\de{Der Fehler kann behoben werden.}}
\exitem{Przetłumacz: „Dane nie mogą być usunięte bez zgody.”}{\de{Die Daten dürfen nicht ohne Genehmigung gelöscht werden.}}
\exitem{Przetłumacz: „Problem da się powtórzyć.”}{\de{Das Problem lässt sich reproduzieren.}}
\end{cwiczenia}

\section{Trening słownictwa (generowany)}
\drill{40}

\weekend{Opisz status swojego projektu w ośmiu zdaniach z użyciem: dwóch \de{werden}-Passiv, dwóch \de{sein}-Passiv, jednego z modalnym
i jednego z \de{sich lassen}. Nagraj 2 minuty jako raport dla kierownika.}
